\section*{Chapter \ref{ch:b2:social-choice} --- Social Choice}

\begin{solution}{exo:b2:social-choice:1}
% code: exercises/solutions.py pairwise condorcet_winner
$A$ beats $B$ by 3 to 2 (the two $A \succ B \succ C$ voters and the $C \succ A \succ B$ voter); $B$ beats $C$ by 4 to 1; $C$
beats $A$ by 3 to 2. The majority's preferences cycle, and there is no \omterm{def:b2:social-choice:condorcet-winner}{Condorcet winner}.
\end{solution}

\begin{solution}{exo:b2:social-choice:2}
$y$ is a \omterm{def:b2:social-choice:pareto-improvement}{Pareto improvement} on $x$: the first person gains, the second is as well off. $z$ cannot be compared
with $x$ or with $y$: the first person prefers $z$, the second prefers $x$ and $y$.
\end{solution}

\begin{solution}{exo:b2:social-choice:3}
% ledger: b2:gsm2023-t7
% code: listings/condorcet_random.py
One election of 182, 0.55 per cent. With random rankings the listing finds no \omterm{def:b2:social-choice:condorcet-winner}{Condorcet winner} in 9.2 per cent of
three-candidate elections and more with more candidates. Real voters' rankings are far from random: they
share a structure, such as an ordering of candidates from one side to the other, that makes a \omterm{def:b2:social-choice:condorcet-winner}{Condorcet
winner} likely.
\end{solution}

\begin{solution}{exo:b2:social-choice:4}
A budget that beats the median voter's favourite in a one-to-one vote, or a winner far from the median
favourite. Before concluding, check that preferences over the budget are \omterm{def:b2:social-choice:single-peaked}{single-peaked} (a voter who wants
either no library or a large one is not), that the vote really was over one dimension, that turnout did not
shift the median of those who voted, and how the median favourite was measured.
\end{solution}

\begin{solution}{exo:b2:social-choice:5}
% code: exercises/solutions.py borda
$A$ beats $B$ and $C$ by 3 to 2: it is the \omterm{def:b2:social-choice:condorcet-winner}{Condorcet winner}. Borda: $A$ has $3 \times 2 = 6$, $B$ has $3 \times 1 + 2
\times 2 = 7$, $C$ has $2 \times 1 = 2$; $B$ wins. The Borda count fails to choose a \omterm{def:b2:social-choice:condorcet-winner}{Condorcet winner} when one exists,
because it counts how far apart alternatives are on each ballot, which depends on where the third
alternative is placed. A single profile cannot show a failure of \omterm{def:b2:social-choice:iia}{independence of irrelevant alternatives},
which compares two profiles; exercise 12 supplies the second: when the three move $C$ above $B$, nobody's
comparison of $A$ with $B$ changes, yet society's ranking of $A$ and $B$ reverses.
\end{solution}

\begin{solution}{exo:b2:social-choice:6}
% code: exercises/solutions.py single_peaked
$L \succ M \succ R$ and $M \succ R \succ L$ are \omterm{def:b2:social-choice:single-peaked}{single-peaked}; $R \succ L \succ M$ is not, since after its peak $R$ it jumps to
$L$, past $M$. A \omterm{def:b2:social-choice:single-peaked}{single-peaked} voter's favourite is $L$, $M$ or $R$; if it is $L$ or $R$, $M$ is next to it and
comes second; if it is $M$, $M$ is first. Either way $M$ is not last.
\end{solution}

\begin{solution}{exo:b2:social-choice:7}
% code: exercises/solutions.py median_shift
Adding two voters above the median should raise the outcome. The median moves from 30 to 40 (the fourth of seven
favourites). Two new voters at 5 and 70, one on each side, leave the median at 30.
\end{solution}

\begin{solution}{exo:b2:social-choice:8}
% code: exercises/solutions.py ex8
$1.5 + 3 \times 3.5 = 12$ and $3 \times 3.5 + 1.5 = 12$: each state is on its policy's frontier. Policy $B$ allows $u_2 =
3.5$ with $u_1 = (12 - 3.5)/3 = 2.83 > 1.5$, a \omterm{def:b2:social-choice:pareto-improvement}{Pareto improvement} on $A$: the move from $A$ to $B$ passes. Policy $A$ allows
$u_1 = 3.5$ with $u_2 = (12 - 3.5)/3 = 2.83 > 1.5$, a \omterm{def:b2:social-choice:pareto-improvement}{Pareto improvement} on $B$: the move back passes too.
\end{solution}

\begin{solution}{exo:b2:social-choice:9}
For each town, regress the approved spending per resident on the median resident's income and on her tax share
(what one more dalen of spending costs her), and compare with the same regression using mean income. The
model predicts that the median does better; spending that tracks mean income, or the preferences of
those who set the agenda, rather than the median, counts against it. Towns with different turnout among
income groups allow a further test: the median of those who vote should predict better than the median of
residents.
\end{solution}

\begin{solution}{exo:b2:social-choice:10}
(a) In every intermediate profile each voter places $b$ at the top or the bottom, so by step 1 society places it at
the top or the bottom; it can only jump. (b) $n^*$ has $b$ at the top in II; moving $a$ just above $b$ changes only
her $a$--$b$ comparison. The other voters keep $b$ at their top or bottom, so they can place $a$ and $c$ in either
order without moving $b$ between them. (c) Step 4 needs a pair whose social ranking $n^*$ is known to change on
her own; between I and II, only her ranking changed and $b$ moved from the bottom to the top socially, so the
voter who decides the pair $(a, b)$ must be $n^*$.
\end{solution}

\begin{solution}{exo:b2:social-choice:11}
Let $n$ be even and the middle favourites $x_{(n/2)} \le x_{(n/2+1)}$, and take $x$ in between. For $y < x$, the $n/2$
voters with favourites at or above $x_{(n/2+1)}$ prefer $x$ to $y$: at least half, so $y$ does not win. For $y > x$, the
$n/2$ voters with favourites at or below $x_{(n/2)}$ prefer $x$. So no alternative beats $x$; with exactly half on each
side, a vote can tie.
\end{solution}

\begin{solution}{exo:b2:social-choice:12}
% code: exercises/solutions.py borda
If the three report $A \succ C \succ B$, the scores are $A$ 6, $C$ $3 + 2 = 5$, $B$ $2 \times 2 = 4$: $A$ wins. The two others
could answer by reporting $C \succ B \succ A$, which gives $A$ 6, $B$ 2, $C$ $3 + 4 = 7$: $C$ wins, which they prefer to
$A$. Each side's best report depends on the other's: the rule invites strategy, as the Gibbard--Satterthwaite
theorem says every non-dictatorial rule with three or more alternatives must.
\end{solution}

\begin{solution}{pb:b2:social-choice:1}
% code: exercises/solutions.py problem part4
\textbf{Part I.}
1. 250: $200 \succ 400 \succ 600 \succ 800$; 380: $400 \succ 200 \succ 600 \succ 800$; 450: $400 \succ 600 \succ 200 \succ 800$;
520: $600 \succ 400 \succ 800 \succ 200$; 640 and 690: $600 \succ 800 \succ 400 \succ 200$; 900: $800 \succ 600 \succ 400 \succ 200$.
2. First places: 200 one, 400 two, 600 three, 800 one. 600 wins.
3. 200: $3 + 2 + 1 = 6$; 400: $2 + 3 + 3 + 2 + 1 + 1 + 1 = 13$; 600: $1 + 1 + 2 + 3 + 3 + 3 + 2 = 15$; 800: $1 + 2 + 2 + 3 = 8$.
4. 600 beats 200 by 5 to 2, 400 by 4 to 3, 800 by 6 to 1.
5. Yes, 600.
6. Yes: plurality, Borda and pairwise majority all choose 600.
\textbf{Part II.}
7. Each ranking starts at the option nearest the favourite and then moves outward: for example 450 ranks 400,
then 600, then 200, then 800, each step to an option next to those already ranked.
8. No: after 200 it jumps to 800, past 400 and 600. It could come from a councillor who wants either a small
garden or a full park with facilities, and sees a middle budget as buying neither.
9. Moving away from the favourite in either direction increases the distance, so utility falls on each side
of the peak.
10. The \omterm{def:b2:social-choice:median-voter}{median voter theorem}: the median councillor's favourite option wins every pairwise vote.
\textbf{Part III.}
11. The fourth of the seven favourites: 520.
12. 600, at distance 80 against 120 for 400.
13. See question 4.
14. The favourites become 380, 450, 520, 640, 690, 900, 950: the median is 640, whose nearest option is still 600; 600
remains the \omterm{def:b2:social-choice:condorcet-winner}{Condorcet winner}.
15. Scores: 200 7, 400 13, 600 12, 800 10: 400 wins. Alone, 380 lowers 600 to 14 and 400 still has 13; 450 alone lowers 600
to 13, a tie with 400. Together they change the result.
16. No. The outcome is the median's favourite option. A councillor below the median who reports a higher favourite
either leaves the median unchanged or moves it up, away from her; one above it can only move it down, also away.
\textbf{Part IV.}
17. $X$ beats $Y$ by 5 to 2; $Y$ beats $Z$ by 5 to 2; $Z$ beats $X$ by 4 to 3.
18. $X \succ Y$, $Y \succ Z$ and $Z \succ X$ for the majority: a cycle.
19. $X$ beats $Y$, then $Z$ beats $X$: $Z$ is adopted.
20. $Y$ against $Z$, then $X$: $Y$ wins, then $X$ wins. $X$ against $Z$, then $Y$: $Z$ wins, then $Y$ wins. Each proposal wins
under some agenda.
21. The rankings are not \omterm{def:b2:social-choice:single-peaked}{single-peaked} on any line through $X$, $Y$ and $Z$: two questions, budget and site, give no
single order along which every councillor's preference falls on both sides of a peak, so the restriction that
rescued Part III fails and the unrestricted domain returns.
22. Yes: $Y$ would beat $X$ by 5 to 2, then $Y$ would beat $Z$ by 5 to 2. They get $Y$ instead of $Z$, which they prefer.
23. Whoever sets the agenda can choose the outcome, as long as the others vote sincerely; councillors who
anticipate the agenda can partly undo that power by voting strategically.
24. The budget of 600 thousand dalens wins every pairwise vote. Once the site is added, the agendas ($X$ v $Y$, then
$Z$), ($Y$ v $Z$, then $X$) and ($X$ v $Z$, then $Y$) adopt $Z$, $X$ and $Y$.
\end{solution}
